% !TeX spellcheck = pt_BR
% ============================================================
% Convencoes e notacao usadas no livro
% ============================================================

\chapter*{Convenções e Notação}
\addcontentsline{toc}{chapter}{Convenções e Notação}

Este capítulo serve de referência rápida para as convenções tipográficas
e notacionais adotadas ao longo do livro.

\section*{Classificação das equações}

Toda equação apresentada é classificada por dois critérios independentes,
indicados visualmente na sua caixa de apresentação:

\begin{itemize}
  \item \textbf{Nível} (I, II ou III): indica a profundidade conceitual.
  \begin{itemize}
    \item \textbf{Nível I:} fundamentais, indicadas a todos os leitores.
    \item \textbf{Nível II:} aprofundamentos para vestibulares concorridos
    (medicina, exatas) ou estudos posteriores.
    \item \textbf{Nível III:} avançadas, voltadas a olímpiadas
    (OBF, IpHO) e vestibulares militares (IME, ITA).
  \end{itemize}

  \item \textbf{Tipo} epistemológico:
  \begin{itemize}
    \item \textbf{Definição} ({def}): cria um conceito novo. Não pode
    ser errada --- é convenção.
    \item \textbf{Experimental} ({exp}): nasce de medições. Síntese
    matemática de como a natureza se comporta.
    \item \textbf{Demonstrável} ({dem}): consequência lógica de
    definições e equações experimentais já estabelecidas.
  \end{itemize}
\end{itemize}

\section*{Notação matemática}

\begin{center}
\renewcommand{\arraystretch}{1.4}
\begin{tabular}{>{\centering\arraybackslash}p{3cm} l}
\toprule
\textbf{Símbolo} & \textbf{Significado} \\
\midrule
$:=$           & Definição (lado esquerdo é definido pelo direito) \\
$\implies$     & Implica (segue logicamente de) \\
$\Delta x$     & Variação da grandeza $x$ ($x_{\text{final}} - x_{\text{inicial}}$) \\
$\vec{v}$      & Grandeza vetorial \\
$v$            & Módulo (ou componente escalar) da grandeza vetorial \\
$\hat{x}$      & Vetor unitário na direção do eixo $x$ \\
$|\vec{v}|$    & Norma (módulo) de um vetor \\
\bottomrule
\end{tabular}
\end{center}

\section*{Notação de unidades}

As unidades adotadas seguem o Sistema Internacional (SI) ao longo de
todo o livro, salvo indicação explícita em contrário. Quando aplicável,
unidades são expressas usando o pacote \texttt{siunitx}, por exemplo:
$g = 9{,}8 \text{ m/s}^2$.

\clearpage
